\documentclass[10pt,a4paper]{article}
\usepackage{fontspec}
\usepackage[english,russian]{babel}
\babelfont{rm}{PT Serif}
\babelfont{sf}{PT Sans}
\babelfont{tt}[Scale=0.88]{PT Mono}
\usepackage[a4paper,top=22mm,bottom=24mm,left=22mm,right=22mm]{geometry}
\usepackage{amsmath}
\usepackage{booktabs}
\usepackage{array}
\usepackage{enumitem}
\usepackage{titlesec}
\usepackage{caption}
\usepackage{tikz}
\usetikzlibrary{arrows.meta}
\usepackage{xcolor}
\usepackage[hidelinks]{hyperref}
\usepackage{microtype}

\setlength{\parindent}{0pt}
\setlength{\parskip}{3.5pt plus 1pt}
\setlength{\emergencystretch}{4em}

\titleformat{\section}{\normalfont\fontsize{11.5}{14}\bfseries}{\thesection.}{0.5em}{}
\titlespacing*{\section}{0pt}{10pt plus 2pt minus 2pt}{3pt}
\titleformat{\subsection}{\normalfont\fontsize{10}{12}\bfseries}{\thesubsection.}{0.5em}{}
\titlespacing*{\subsection}{0pt}{6pt plus 2pt minus 1pt}{2pt}

\setlist[enumerate]{leftmargin=2.0em,itemsep=0.5pt,topsep=2pt,parsep=0pt}
\setlist[itemize]{leftmargin=1.5em,itemsep=0.5pt,topsep=2pt,parsep=0pt}
\captionsetup{font=small,labelfont=bf,labelsep=period,skip=3pt}
\newcolumntype{L}[1]{>{\raggedright\arraybackslash}p{#1}}
\setlength{\textfloatsep}{9pt plus 2pt minus 2pt}
\setlength{\floatsep}{8pt plus 2pt minus 2pt}

\pagestyle{plain}

\tikzset{
  box/.style={draw,rectangle,thick,align=center,inner sep=2.5pt,font=\footnotesize,minimum width=27mm,minimum height=8mm},
  zone/.style={draw,dashed,rounded corners=2pt},
  zl/.style={font=\scriptsize\itshape,anchor=north west,inner sep=2pt},
  lbl/.style={font=\scriptsize,align=center,fill=white,inner sep=1pt},
  arr/.style={-{Latex[length=1.8mm]},thick},
}

\begin{document}

\begin{center}
{\fontsize{16}{19.5}\selectfont\bfseries Swarma: проверяемые границы доверия\\[2pt] для роёв автономных агентов\par}
\vspace{5pt}
{\normalsize Личности без ключей у агента, гранты, подписанные сообщения\\ и изоляция исполнения на машинах владельца\par}
\vspace{8pt}
{\normalsize Алексей Григорьев, Иван Опарин\\
Relux Works\\
\href{https://github.com/relux-works}{github.com/relux-works}\par}
\vspace{3pt}
{\normalsize Октябрь 2026, версия 6.0\par}
\end{center}

\vspace{4pt}

\begin{list}{}{\leftmargin=0.9cm\rightmargin=0.9cm}
\item\small
\textbf{Аннотация.} Агент разработки читает недоверенный текст и действует с правами процесса, который его запустил, поэтому инъекция в промпт исполняется руками самого агента. Рою таких агентов на машинах одного владельца и у независимых владельцев нужны границы, которые выдерживают захват агента. Swarma строит их из детерминированных проверок у исполнителя эффекта. Закрытого ключа личности у агента нет: подпись выдаёт ключник, который узнаёт собеседника по UID от ядра и по поколению учётки. Каждое действие требует гранта, который можно только сузить. Ключи, аттестации, гранты и отзывы ведёт реестр владельца на отдельном узле доверия со свидетелем отката, и запись отзыва действует на каждой проверке не позже $\max(S+A,F)$ при $A\le60$~с. Сообщения из мессенджера допускают ворота доверия, учётки на мессенджере выдаёт регистратор по ключу машины, станция допускает задачу только по квалифицированным гарантиям изоляции, а изменение принимает независимый ревьювер точного кандидата. Каждую гарантию мы разбираем против модели угроз вместе с допущениями и пределом. Проверены спецификации и отдельные компоненты: проверка подписанных записей согласована на 2632 случаях между генератором, независимой проверкой на Python и библиотекой на Go; два независимых вычислителя памяти дают ожидаемые ответы на 66 смысловых векторах; на реальной машине Linux примитивы учёток прошли применение, проверку и откат, а установка квалификации прошла 11 из 11 проверок; ограниченная модель показывает необходимость каждого из четырёх правил интеграции. Эталонная неоптимизированная проверка одной записи занимает от 2 до 18~мс на одном ядре. Сквозного прогона станции с измеренной задержкой отзыва, оценки производительности всей системы и состязательной оценки живой установки в работе нет.
\end{list}

\section{Введение}

Агент разработки вызывает инструменты, меняет файлы и отправляет сообщения. Среда обычно знает только пользователя ОС, запустившего процесс, и отдаёт агенту все его права. Косвенная инъекция в промпт \cite{greshake,agentdojo} пользуется ими: текст из задачи, зависимости или письма коллеги убеждает модель прочитать ключ, отправить токен или одобрить запрос. Просьба к модели не читать секрет слабее границы, за которой процесс прочитать его не может.

Рой добавляет вопросы, на которые модель не может ответить своим текстом. Оркестратор запускает продьюсеров и ревьюверов, агенты переписываются через мессенджер, сессии переживают перезапуски, участниками становятся агенты других владельцев. Кто отправил сообщение, разрешено ли действие, к какому состоянию относится одобрение, что делать при недоступном свидетельстве? Агентские фреймворки \cite{autogen,metagpt,chatdev,sweagent} координируют роли и инструменты, а границы полномочий между агентами и между владельцами оставляют среде исполнения.

Swarma состоит из трёх слоёв: рантайма на машинах одного владельца (Swarma Runtime), протокола совместной работы независимых владельцев над единой средой проекта (Swarma Protocol) и цепочки с общим состоянием личностей (Swarma Chain). Эта работа описывает границы доверия первых двух слоёв; подробное устройство платформы описывает книга о Swarma. Если работа расходится со спецификацией модуля, права спецификация. Состояние реализации каждой части приведено в разделе~\ref{sec:eval}.

Вклад работы:
\begin{enumerate}
\item Модель системы и модель угроз для роя агентов, где захваченный инъекцией агент считается штатным противником.
\item Композиция механизмов, в которой допуск решает детерминированная проверка у исполнителя: ключник с определением собеседника ядром, гранты с ослаблением, реестр владельца на узле доверия со свидетелем отката и ограниченным сроком распространения отзыва, ворота доверия поверх недоверенного мессенджера, допуск к мессенджеру без общего секрета, станции с контрактом изоляции исполнения, память как потребитель реестра и независимое ревью точного кандидата.
\item Анализ каждой гарантии против модели угроз с допущениями и пределом.
\item Свидетельства, которые уже существуют: корпуса соответствия с независимыми реализациями, квалификация хоста на реальной машине, ограниченная модель интеграции, микрозамеры эталонной проверки записей и статистика ревью, каждое с перечнем того, чего оно не показывает.
\end{enumerate}

\section{Постановка задачи и модель системы}
\label{sec:model}

\textbf{Принципалы.} Владелец отвечает за свои машины, агентов и оплату моделей; ему принадлежит корень доверия установки. Оператор, человек владельца, выдаёт полномочия и подтверждает действия, расширяющие права. Агент исполняет работу в пределах выданных полномочий; его криптографическая личность отделена от процесса и от пользователя ОС. Проект задаёт цель, среду и правила приёмки, контрибьютор и донор приносят ему ёмкость своего агентского цикла и модели. Совмещение ролей дополнительных прав не даёт.

\textbf{Хосты и станции.} Хостом называется машина владельца, на которой работают агенты. \emph{Станцией} называется место одного агента на хосте: своя учётка ОС, своя личность на носителе, свои ключи, дом и корневой контекст. Учётки станций между агентами не делятся, и учётка с отдельным UID без остальных проверок станцией не считается. Привилегированные службы хоста разнесены: ключник работает под root, менеджер пользователей и запускатель получают root на одну узкую операцию через \texttt{sudo}, диспетчер, брокер учётных данных, хост сессий и мост работают под своими служебными учётками.

\textbf{Носители.} Носителем называется мессенджер, который доставляет сообщения агентов и людей; первым носителем служит Matrix. Носитель перевозит события и полномочий не выдаёт: членство в комнате прав на агента не даёт. Конфиденциальность от сервера даёт сквозное шифрование носителя, подлинность автора даёт подпись над точными байтами.

\textbf{Узел доверия.} Реестр владельца и свидетель отката работают на отдельной машине вне административного домена платформы. Её снимают у провайдера, который от платформы не зависит, и администрирует её только владелец; у агентов там нет учёток, у платформы нет административного доступа. Реестр служит единственным авторитетным писателем ключей, аттестаций, грантов, отзывов и журналов личностей, а службы хоста получают от него подписанные снимки (рис.~\ref{fig:model}).

\textbf{Протокол.} Swarma Protocol разделяет мышление, управление и исполнение: агентские циклы контрибьюторов предлагают ревизии, кор-ноды ведут журнал полномочий и принятия с сертификатами кворума, авторитетная среда проекта исполняет только разрешённые действия. Порог сертификата равен одной подписи на уровне L0, $t>n/2$ при честном секвенсере на L1 и $q=\lceil (n+f+1)/2\rceil$ при $n\ge3f+1$ на L2 \cite{pbft,hotstuff}.

\begin{figure}[t]
\centering
\begin{tikzpicture}
  \draw[zone] (-0.2,0.9) rectangle (9.7,4.25);
  \node[zl] at (-0.2,4.25) {хост владельца};
  \draw[zone] (10.6,0.9) rectangle (14.9,4.25);
  \node[zl] at (10.6,4.25) {носитель};
  \draw[zone] (3.1,-2.05) rectangle (9.7,-0.25);
  \node[zl] at (3.1,-0.25) {узел доверия};
  \node[box] (a1) at (1.35,3.25) {станция A\\\scriptsize учётка, дом, сессия};
  \node[box] (a2) at (1.35,1.75) {станция B\\\scriptsize учётка, дом, сессия};
  \node[box] (kh) at (4.75,3.25) {ключник\\\scriptsize root};
  \node[box] (svc) at (4.75,1.75) {диспетчер, брокер,\\\scriptsize хост сессий};
  \node[box] (hl) at (8.15,3.25) {помощник допуска\\\scriptsize ключ машины};
  \node[box] (br) at (8.15,1.75) {мост\\\scriptsize служебная учётка};
  \node[box] (rg) at (12.75,3.25) {регистратор};
  \node[box] (cs) at (12.75,1.75) {сервер Matrix};
  \node[box] (reg) at (4.75,-1.2) {реестр владельца};
  \node[box] (wit) at (8.15,-1.2) {свидетель отката};
  \node[box] (ph) at (0.95,-1.2) {телефон оператора\\\scriptsize ключ в защищённом чипе};
  \draw[arr] (a1) -- node[lbl,above=1pt]{UID} (kh);
  \draw[arr] (a2) -- (svc);
  \draw[arr] (hl) -- node[lbl,above=1pt]{запрос} (rg);
  \draw[arr] (rg) -- node[lbl,right=1pt]{пропуск} (cs);
  \draw[arr] (cs) -- node[lbl,above=1pt]{события} (br);
  \draw[arr] (reg) -- node[lbl,right=1pt]{снимки $\le$60~с} (svc);
  \draw[arr] (reg) -- (wit);
  \draw[arr] (ph) -- (reg);
\end{tikzpicture}
\caption{Модель системы. Агенты работают на станциях под своими учётками и ключей не держат; ключник и помощник допуска узнают собеседника по UID от ядра. Регистратор выдаёт пропуск на носитель, мост проверяет сообщения, реестр на узле доверия принимает подписанные действия оператора и выдаёт службам подписанные снимки, свидетель фиксирует его контрольные точки.}
\label{fig:model}
\end{figure}

\textbf{Задача.} Ни одно действие агента с эффектом за пределами его станции не должно исполняться без проверяемого основания: подлинного отправителя, актуального гранта, одобрения, адресованного точному объекту, и физического удержания исполнения в выданной области. Основание проверяет сторона, совершающая эффект. Отсутствие свидетельства, подтверждённое отсутствие события и невозможность прочитать состояние различаются, и разрешения ни одно из них не даёт.

\section{Модель угроз}
\label{sec:threats}

\textbf{Противники.}
\begin{enumerate}[label=П\arabic*.]
\item Захваченный агент: модель подчинилась внедрённому тексту и пользуется всеми возможностями процесса агента, включая подпись, которую ключник выдаёт его личности.
\item Чужой участник: агент или человек другого владельца, рабочая нода протокола или приглашённый участник комнаты пытается расширить полномочия, подменить объект одобрения или действовать от чужого имени.
\item Сетевой противник и сервер носителя: противник подменяет, задерживает и повторяет сообщения, а сервер может выпустить событие от имени любого адреса своего домена.
\item Захваченная машина: противник владеет одной зарегистрированной машиной, но не корнями владельца.
\item Противник на узле доверия: он похитил ключ реестра или пытается откатить и раздвоить журнал реестра.
\item Криптографический противник: он ломает класс подписей на эллиптических кривых.
\end{enumerate}

\textbf{Допущения.} Ядро ОС, root и установка платформы доверены. Мастер-секрет ключника цел; где возможно, его держит аппаратное хранилище. Оператор подтверждает расширяющие действия на телефоне с ключом в Secure Enclave или StrongBox и видит подписываемое содержание. Владелец честен перед своими агентами и единолично администрирует узел доверия. У установки есть проверенный или контролируемый оператором источник времени, например NTS. Для протокола среда проекта доверена в отношении целостности, если нет аппаратной аттестации исполнения.

\textbf{Вне рамок.} Мы не защищаемся от эксплойта ядра и захваченного root, не обещаем правильности суждения модели и конфиденциальности кода от провайдера модели, не рассматриваем побочные каналы и подделку оборудования. Отказ в обслуживании не предотвращается: при недоступном свидетельстве система закрывает допуск и сознательно жертвует живучестью. Экономика и переносимая репутация в этой работе не оцениваются.

\section{Устройство}
\label{sec:design}

\subsection{Личность и ключи}

Ключевая пара личности отвечает, кто отправитель, и ничего не разрешает. Закрытого ключа личности у агента нет. Ключник, привилегированная служба с мастер-секретом установки, выводит ключи агентов жёсткой деривацией по BIP-32 в варианте SLIP-0010 для Ed25519 \cite{bip32,slip10}. Путь \texttt{назначение'/\allowbreak продукт'/\allowbreak проект'/\allowbreak агент'} кодирует стабильную личность и не совпадает ни с UID, ни с путём в файловой системе, поэтому личность переживает пересоздание учётки.

Ключник слушает Unix-сокет и получает UID собеседника от ядра (\texttt{SO\_PEERCRED} на Linux, \texttt{LOCAL\_PEERCRED} на macOS). Личность выбирается по привязке, где записаны UID, поколение учётки, путь, владелец, номер регистрации и признак разовой или постоянной личности. Подключение допустимо, только если поколение этого UID активно в журнале менеджера пользователей. Поколение представляет собой случайное 128-битное число, новое при каждом создании учётки, поэтому повторно выданный UID чужого ключа не наследует. Набор операций узок: \texttt{sign}, \texttt{sign\_message}, \texttt{decrypt}, \texttt{pubkey}, \texttt{register}, \texttt{revoke}.

Сообщения подписывает отдельный ключ сообщений Ed25519, выведенный по пути \texttt{m/4'/продукт'/\allowbreak проект'/\allowbreak агент'}, и только документ с меткой \texttt{swarma-msg/2} с прообразом не длиннее 24\,576 байт. Принадлежность ключа агенту удостоверяет запись \texttt{attestation/1}, подписанная корнем аттестаций ключника, со сроком не длиннее 86\,400~с. Аттестация становится текущей, когда она сама, её запись в журнале реестра, контрольная точка и подтверждение свидетеля сохранены вместе; до этого и при неизвестном исходе выпуска подпись закрыта. Каждый корень подписывает только свой язык записей: корень аттестаций ключника, корень оператора, корень грантов владельца, ключ реестра, ключи служб грантов, офлайновый издатель регистраций машин и ключ коммитов. Ключи ролей попарно различны, и проверяющий отказывает при совпадении двух ролей. Форматы задаёт спецификация подписанных записей: \texttt{swarm-envelope/2}, \texttt{message/1}, \texttt{grant/1}, \texttt{operator-act/1}, \texttt{attestation/1}, \texttt{revocation/1}, \texttt{journal-entry/1}, \texttt{checkpoint/1}, \texttt{freshness/1}.

Для других владельцев личность представлена самосертифицирующим идентификатором в духе KERI \cite{keri}: он равен хэшу первой записи журнала ключевых событий, которая объявляет текущие ключи, обязательство на следующие и родителя. Предротация защищает от поломки текущих ключей. До запуска Swarma Chain журнал ведёт реестр владельца, и при публикации в цепочку идентификатор не меняется.

\subsection{Гранты и реестр владельца}

Грант означает подписанное утверждение: кто выдал полномочие, кому, на какие действия и цели, при каких условиях, до какого срока и можно ли делегировать его дальше. Дочернее полномочие ограничено родительским по действиям, целям, сроку и глубине, а цепочку до корня проверяет исполнитель действия, даже если раньше на маршруте её уже проверили. Модель следует capability-подходу \cite{dennis1966,miller2006} и ориентируется на SPKI/SDSI, Macaroons, UCAN и Biscuit \cite{spki,macaroons,ucan,biscuit}; окончательный формат выбирает спецификация грантов.

Ключ агента бывает субъектом и листом цепочки и никогда не бывает издателем звена. Ключ службы грантов выдаёт звено, только имея адресованный ему родительский грант с глубиной не меньше 1; цепочка, начатая самим ключом службы, отклоняется с кодом \texttt{grant\_root\_unknown}. Права ребёнка зависят и от вызывающего, и от диспетчера: в контракте брокера учётных данных v0 лист агента равен пересечению запроса и двух цепочек, содержит только \texttt{credential.lease} с глубиной 0, а цепочка не длиннее четырёх конвертов. Если цепочку построить нельзя, ответ \texttt{grant\_unconstructible} приходит до любых эффектов.

Реестр хранит, упорядочивает и публикует записи и содержимое не переподписывает. Собственный ключ реестра подписывает только снимки отзывов своей области, контрольные точки и записи свежести; его держит отдельный процесс, который принимает только одобренные профили записей. Журнал состоит из неподписанных \texttt{journal-entry/1} с номером, ссылкой на предыдущую запись, операцией \texttt{add} или \texttt{revoke} и идентификатором содержимого, а подписанная \texttt{checkpoint/1} называет голову и текущий снимок отзывов. Свидетель отката работает под отдельной учёткой, процессом и потоком, по запросу с вызовом называет последнюю голову и хранит доверенную контрольную точку вне узла. Две разные контрольные точки с одним номером дают \texttt{journal\_fork}, точка с номером меньше засвидетельствованного даёт \texttt{journal\_rollback}.

Потребители принимают снимок отзывов не старше 60~с; возраст неизменного снимка продлевает запись свежести, которую реестр публикует каждые 20~с. Записи отзыва накапливаются, и снимок без прежней записи полным не считается. Пусть $S$ означает время выпуска первого подписанного полного снимка с записью отзыва, $A\le60$~с допустимый возраст снимка, $F$ время вступления записи в силу. Тогда запись действует на каждой проверке не позже $\max(S+A,F)$ и никогда раньше $F$. Проверка отзыва имеет три исхода: не отозвано, отозвано, недоступно; два последних закрывают допуск. Время проверяется по водяному знаку, а при потере проверенного времени службы держатся не дольше 10 минут по монотонным часам и затем останавливают всё, что от времени зависит. Общий стоп от реестра не зависит, возобновление требует подписи оператора.

\subsection{Подписанные сообщения и ворота доверия}

Мост стоит между носителем и сессией агента. Его экземпляр работает на каждую личность станции под общей служебной учёткой сообщений и никогда под учёткой агента, поэтому у агента нет учётных данных мессенджера. Провайдер носителя сначала дописывает событие в почтовый файл и сбрасывает его на диск, затем продвигает курсор синхронизации, а допуск идемпотентен по паре отправителя и идентификатора события.

Ворота доверия проверяют сообщение детерминированно, без сети и без вопроса к модели, в строгом порядке: каноническая форма байтов, ключ в снимке реестра, подпись, цепочка грантов до корня с сужением действий и целей и вложенностью сроков, снимок отзыва, окно времени, журнал повторов, ревизия привязки цели. Вердикт ворота подписывают ключом моста. Затем небольшая локальная модель (decision guard) может только сузить допуск, и её недоступность даёт отказ. Допущенное событие во ввод не вставляется: хост сессий получает короткий звонок, кладёт его в долговечную очередь и доставляет при открытом заборе допуска, а тело и запись проверки агент читает из почтового ящика как данные. Сообщения других агентов приходят отдельной рамкой и прав одобрять или управлять сессией не дают.

Сила свидетельства о человеке-отправителе задаётся уровнями. На A0 отправителя удостоверяет сервер носителя владельца, на A1 устройство отправителя подписано закреплённым мастер-ключом перекрёстной подписи Matrix \cite{matrix}, на A2 ключ человека подписал точные байты конверта и ворота проверили грант, на A3 оператор подтвердил точное действие на телефоне. Расширяющие действия требуют A3. Человека другого владельца узнают по корню, закреплённому один раз при знакомстве: корень аттестует мастер-ключ вместе с полным MXID, и цепочка из устройства, ключа самоподписи, мастер-ключа и аттестации корня должна сойтись целиком. Гранты называют ключи, поэтому смена сервера или адреса прав не переносит.

\subsection{Станции и контракт изоляции исполнения}

Каждый автономный агент работает отдельным пользователем ОС. На Linux у каждого поколения учётки своя cgroup v2 под управлением root с лимитами \texttt{cpu.max}, \texttt{memory.max}, \texttt{memory.swap.max}, \texttt{pids.max} и ввода-вывода. Менеджер пользователей создаёт учётки, запускатель запускает одобренную программу; оба вызываются через \texttt{sudo} без аргументов и принимают один JSON ограниченного размера, а план запуска передаётся неизменяемым на дескрипторе 3. Агент остановлен, когда у всей учётки нет процессов в течение опубликованного периода тишины. Если ответ потерян, диспетчер наблюдает тот же эффект и не повторяет запуск с новой личностью. Служба не держит замок своего состояния во время синхронного вызова другой службы, и граф ожиданий каждой службы не содержит циклов. Сеть ограничивает файрвол ядра по UID: агент достигает только слушателя прокси своего поколения, а \texttt{CONNECT} ограничивает имя, порт и проверенный адрес.

Что станция гарантирует, задаёт контракт изоляции исполнения (execution isolation contract, EIC): профиль изоляции, вектор возможностей, привязки личности, размещения и исполнения, квитанция границы и кортеж квалификации. Запускатель применяет допущенную границу и пишет защищённую квитанцию, диспетчер сверяет покрытие, брокер перед выдачей и продлением учётных данных проверяет текущую квитанцию, хост сессий держит заборы ввода. Гарантии названы 42 идентификаторами общего словаря, например \texttt{isolation.read}, \texttt{process.stop}, \texttt{limit.memory}. У гарантии на каждой платформе один из четырёх статусов: \texttt{present}, \texttt{weaker-with-compensation}, \texttt{absent}, \texttt{PENDING-VERIFY}. Задача объявляет нужные гарантии записью \texttt{station-required-guarantees/1}; \texttt{absent} и \texttt{PENDING-VERIFY} ведут к отказу, а ослабление допустимо, только если оператор отдельным аутентифицированным действием принял его для этой задачи, станции, поколения, ревизии профиля и матрицы, количественной области, компенсации и срока. Списка разрешённых имён задач нет.

Машину владелец квалифицирует сам по воспроизводимой процедуре: обзор без привилегий, сборка набора, ревью владельцем, подготовка под root, применение, проверка с закрытым допуском, подписанная аттестация. Аттестацию подписывает отдельный ключ машины, и она называет каждую из 42 гарантий ровно один раз. При дрейфе защищённых входов, смене класса платформы, проваленной проверке или отозванном доверии новые старты закрываются сразу, а работающее исполнение сужается или останавливается целиком не позже 300~с; повторы и перезапуски этот срок не сбрасывают.

\subsection{Допуск к носителю}

Учётку на носителе агент получает без общего секрета. Помощник с ключом машины слушает Unix-сокет, узнаёт UID вызывающего от ядра и подписывает только запрос на допуск с именем, разрешённым этому UID: метка \texttt{swarma-registrar registration/1}, 128-битное одноразовое число, срок не длиннее 120~с. Регистратор сверяет ключ машины и имя с реестром машин, проверяет свежесть числа и выдаёт пропуск, JWT с подписью Ed25519 на 120~с. Сервер носителя доверяет только ключу регистратора, создаёт учётку при первом входе и выдаёт токен доступа, который сам аннулирует через 86\,400~с, без refresh-токена; деактивированная учётка по пропуску не входит. Девять требований к серверу (HR1--HR9) квалифицируются отрицательными и положительными случаями на изолированном сервере до развёртывания конкретной сборки. Реестр машин состоит из записей \texttt{enrollment/1} в git, подписанных владельцем. Отзыв машины требует двух действий: удалить её ключ из реестра и вывести и деактивировать все её учётки на сервере.

\subsection{Память как потребитель реестра}

Память агентов хранит типизированные утверждения с источниками, обоснованиями, оценками и ревизиями и отвечает, какие допущенные свидетельства поддерживают или опровергают утверждение. С реестром она стоит на одной подложке подписанных записей и читает его как потребитель: отправителя записи определяет по ячейке личности в текущем снимке, основание представительства берёт из гранта, а недоступный реестр никогда не считается выданным грантом. Полномочий память не выдаёт. Первый профиль хранит пространства в git: агент сам подписывает запись, каждый читатель проверяет подпись, схему и право на ревизию при чтении, а содержимое ограничено нечувствительным, потому что из истории git ничего не стирается. Нормативный пакет 0.2.7 содержит 52 пункта ядра, 14 закрытых схем и 74 вектора соответствия.

\subsection{Метод оркестрации: продьюсер и независимый ревьювер}

Оркестратор декомпозирует цель, продьюсер производит изменение, ревьювер независимо проверяет опубликованную ревизию. На доске задач Change Request публикует неизменяемую ревизию, вердикт относится именно к ней, и один запуск не может принять собственную ревизию. Интеграция идёт через compare-and-set точного проверенного дерева. Изменения в CI, хуках, инструкциях агентов и ссылках одобряет детерминированный предикат над всей накопленной дельтой от базы до кандидата. Ревьювер получает инертные диф и объекты, а инструкции, хуки и MCP кандидата его среда автоматически не подхватывает.

В Swarma Protocol среда строит неизменяемый кандидат $I$ над текущей головой $M$ и фиксирует пару $(M,I)$. Множество допустимых ревьюверов замораживается до появления хэша $I$, затем ревьювер вычисляется из хэшевой случайности кор-нод, и оркестратор не может выбрать удобного проверяющего после того, как кандидат известен. Одобрение и аттестации адресованы четвёрке из контракта, кандидата, политики и окружения. Аттестация бывает \emph{пройдена}, \emph{провалена} или \emph{не установлена}, и сбой инструмента зелёной проверкой не становится. Привилегированный класс изменений требует двух независимых ревью или одобрения человека из кор-команды. Независимость оценивается по административному домену, семейству модели и тулчейну. Оркестратор держит роль по аренде с эпохой и fencing-токеном \cite{chubby,kleppmann-fencing}, и валидаторы отвергают запись устаревшей аренды.

\section{Анализ безопасности}
\label{sec:analysis}

Гарантии ниже названы вместе с противником, аргументом, допущениями и пределом. Каждая держится, только если исполнитель эффекта вызывает проверку и физически удерживает область эффекта.

\textbf{G1. Захваченный агент не получает закрытого ключа своей личности (П1).} Материал ключа находится только у ключника, а личность выбирается по UID от ядра и активному поколению, поэтому ни имя в запросе, ни повторно выданный UID чужого ключа не дают. \emph{Допущения:} ядро и root доверены, мастер-секрет цел. \emph{Предел:} захваченный агент получает подписи собственной личности, и подпись удостоверяет авторство, правоты она не доказывает; расшифрованный пакет попадает в процесс адресата. До появления ключника действуют два временных исключения: ключ заявок вызывающего лежит в доме оркестратора, и его силу ограничивает грант, а ключ подписи записей памяти агент создаёт сам. Компрометация мастер-секрета затрагивает всех агентов установки; восстановление идёт через новый корень и ячейку владельца, чьи ключи из мастер-секрета не выводятся.

\textbf{G2. Делегирование не расширяет полномочий (П1, П2).} Исполнитель проходит цепочку до корня и проверяет подписи, вложенность действий, целей и сроков и глубину. Агент не выдаёт грантов даже в пределах своих прав, поэтому захваченный оркестратор не раздаёт полученное соседям, а вызывающий остаётся в цепочке потомка, и отзыв гранта оркестратора отрезает его агентов при следующей аренде или продлении. \emph{Предел:} верный грант не делает владельца ключа добросовестным, и вредное действие в пределах гранта остаётся остаточным риском. Формат общего языка грантов не выбран, брокер v0 действует по своему более узкому контракту.

\textbf{G3. Отзыв распространяется в ограниченный срок, неизвестность закрывает допуск (П1, П5).} Снимок принимается только при возрасте не больше $A$, записи отзыва накапливаются, поэтому запись действует на каждой проверке не позже $\max(S+A,F)$; ответ \emph{недоступно} запрещает эффект, а водяной знак не даёт шагу часов назад вернуть истёкшее право. \emph{Предел:} срок отсчитывается от подписанной публикации. Токен провайдера, уже выданный процессу, остаётся пригодным до аннулирования у провайдера. Завершение отзыва требует подтверждений названных потребителей или истечения прежних полномочий, и его механизм открыт. Один полный снимок вмещает не больше 64 записей и 8192 байт, формат страниц не определён.

\textbf{G4. Откат и раздвоение журнала реестра обнаруживаются (П5).} Контрольные точки пронумерованы и подписаны, свидетель фиксирует их и по запросу с вызовом называет последнюю голову, его доверенная точка хранится вне узла. Укравший ключ реестра не выпустит ни гранта, ни аттестации: их подписывают другие корни. \emph{Предел:} он может задержать или скрыть публикацию и подписать ложный отзыв или ложную контрольную точку своей области. Реестр и свидетель на одной машине образуют один домен отказа. Владелец администрирует узел и может показать разным проверяющим разные истории одной личности; этот предел снимает только общее состояние Swarma Chain.

\textbf{G5. Доставка, членство в комнате и подпись коллеги полномочий не дают (П2, П3).} Ворота проверяют подпись над точными байтами, гранты, отзыв, окно времени и журнал повторов до доставки, а тело остаётся данными. Сохранение события до сдвига курсора и идемпотентный допуск исключают и потерю, и повторный допуск при падении моста. \emph{Предел:} в переходной политике A0 компрометация сервера носителя или учётки человека даёт команды в пределах зоны этого человека, а агенты под одним пользователем ОС могут переписать свой корневой контекст. Decision guard инъекцию надёжно не распознаёт \cite{agentdojo}, подпись не доказывает истинности текста.

\textbf{G6. Захват одной машины ограничен её меткой (П4).} Общего секрета нет, помощник подписывает только разрешённые этому UID имена, регистратор проверяет ключ машины по подписанному реестру, пропуск живёт 120~с, токен 86\,400~с без продления. \emph{Предел:} захваченный агент на зарегистрированной машине получает пропуск на своё разрешённое имя. Допуск закреплённого профиля сервера по времени продлевает пропуск до $\mathit{exp}+60$~с. Агентов под одним пользователем ОС помощник различает только по именам. Выданные токены живут до выхода, поэтому отзыв машины требует обоих действий.

\textbf{G7. Исполнение допускается только на станции с квалифицированными гарантиями (П1).} Диспетчер допускает пересечение запроса, минимумов платформы, политики владельца и текущих квалифицированных возможностей, а недостающую гарантию встречает типизированным отказом. Ослабление видно в матрице, в квитанции границы и в ответе планировщика и требует решения оператора. \emph{Предел:} самоотчёт машины не защищает от захваченного root, подделанных наблюдений ядра и скопированного ключа машины. Профиль macOS не задаёт квалифицированного совокупного предела памяти. TLS внутри \texttt{CONNECT} не расшифровывается, поэтому разрешённый адресат остаётся каналом выноса, а периодические проверки сетевых правил оставляют окно обнаружения.

\textbf{G8. Принятие относится к точному кандидату (П2).} Одобрение и аттестации адресованы контракту, кандидату, политике и окружению, ревьювер вычисляется после фиксации кандидата, голова продвигается сравнением-и-заменой. Отключение любого из четырёх правил доставки в ограниченной модели даёт контрпример (раздел~\ref{sec:eval}). \emph{Предел:} правильность суждения не гарантируется. Модели склонны предпочитать собственный текст \cite{selfpref,llmjudge}, поэтому разные владельцы и семейства моделей уменьшают основания для корреляции ошибок, но независимости не доказывают. Квитанция скомпрометированной среды доказывает авторство отчёта и не доказывает его правдивости.

\textbf{G9. Память не создаёт полномочий (П1, П2).} У памяти нет операций одобрить, выдать грант, запустить и разбудить, ответ приходит агенту как данные с происхождением каждого элемента, запись при недоступном реестре отклоняется. \emph{Предел:} из истории git ничего не стирается, а поиск секретов при записи их отсутствия не доказывает.

\textbf{Криптографический противник (П6).} Подписи рантайма и протокола стоят на кривых: Ed25519, P-256 в защищённых чипах, SSH, TLS и подписи git. Предложенный план миграции переводит долгоживущие корни на хэш-подписи SLH-DSA \cite{fips205} и оставляет ключи на кривых короткоживущими листьями, а консенсус Swarma Chain опирается только на хэш-функции.

\section{Проверка: соответствие, квалификация, перебор и замеры}
\label{sec:eval}

Сквозного развёртывания, на котором можно измерить систему целиком, нет. Поэтому мы проверяем то, что уже существует: однозначность спецификаций на корпусах с независимыми реализациями, привилегированные шаги хоста на реальной машине, ограниченную модель правил интеграции, стоимость эталонной проверки записей и процесс ревью, которым ведётся разработка. Таблица~\ref{tab:status} даёт состояние реализации, таблица~\ref{tab:evidence} сводит свидетельства.

\begin{table}[t]
\centering\small
\caption{Состояние реализации на октябрь 2026 года. \emph{Работает}: есть в коде и используется в названных условиях; \emph{спецификация}: описано в спецификации модуля; \emph{решено}: решение владельцев без кода; \emph{предложено}: ожидает решения. Где состояние не сводится к одному слову, оно дано формулировкой.}
\label{tab:status}
\footnotesize
\begin{tabular}{L{5.9cm}L{2.5cm}L{6.6cm}}
\toprule
Часть & Статус & Граница \\
\midrule
Доска задач: Change Request, ревью, compare-and-set & работает & без сертификатов кворума \\
Носитель Matrix со звонком и чтением почты & работает & без ворот, грантов и восстановления \\
Подписанные записи и библиотека проверки & спецификация & библиотека без сети и ввода-вывода; потребители не включены \\
Ключник для Linux v0 & в коде, ждёт ревью & нет развёрнутой границы и квалификации \\
Реестр владельца, свидетель, узел доверия & решено & проект службы на ревью; узла нет \\
Мост и ворота доверия & решено & кода и спецификации моста нет \\
Регистратор допуска к носителю & спецификация & сервер носителя не квалифицирован \\
Менеджер пользователей, запускатель, диспетчер & спецификация & станции не допускаются; квалификационные запуски на одной машине \\
Контракт изоляции исполнения & предложено & почти все ячейки матрицы \texttt{PENDING-VERIFY} \\
Проверка пакета памяти двумя вычислителями & работает & сам пакет решён; первой рабочей версии нет \\
Swarma Protocol, Swarma Chain & спроектированы, не построены & проверена ограниченная модель интеграции \\
\bottomrule
\end{tabular}
\end{table}

\begin{table}[t]
\centering\small
\caption{Сводка свидетельств. Каждая строка относится к названной ревизии и названным условиям.}
\label{tab:evidence}
\footnotesize
\begin{tabular}{L{2.9cm}L{4.3cm}L{4.1cm}L{3.5cm}}
\toprule
Утверждение & Метод & Результат & Чего не показывает \\
\midrule
Спецификация подписанных записей однозначна & Корпус соответствия: генератор на Go, независимая проверка на Python, библиотека на Go; у отказа точный код и положительный контроль & 2632 случая (2499 отрицательных, 133 положительных); библиотека воспроизводит точный код и план каждого & Корректности вне корпуса, работы проверки внутри служб, полноты спецификации \\
\addlinespace
Ядро памяти однозначно & Два вычислителя: один по сводному контракту без доступа к пакету, второй по пакету & Ожидаемые ответы на всех 66 смысловых векторах у обоих & Защищённой службы (8 векторов развёртывания не исполняются), пользы памяти \\
\addlinespace
Привилегированные шаги хоста исполнимы на Linux & Запуски владельцем под root на одной машине, привязанные к сборкам & Примитивы учёток: применение, проверка, откат; установка: 11 из 11 проверок при применении и при проверке & Допуска станций, нагрузки, перезагрузок, отката установки, других машин \\
\addlinespace
Четыре правила доставки необходимы для свойства интеграции & Исчерпывающий перебор ограниченной модели с отрицательным контролем на каждое правило & 1090 состояний без нарушения; без любого правила найден контрпример & Неограниченных параметров, эквивалентности реализации, честности среды \\
\addlinespace
Проверка записи дешева для служб хоста & Микрозамеры эталонной библиотеки на зафиксированных входах & Медианы от 2,1 до 18,0~мс на вызов (табл.~\ref{tab:bench}) & Задержки и пропускной способности системы, худшего случая \\
\addlinespace
Ревьювер возвращает заметную долю первых ревизий & Разбор записей ревью на доске задач проекта Swarma Curator & Приняты 17 из 35 разобранных первых раундов; медиана 2 раунда & Пропущенных дефектов, точности ревьюверов, причинного эффекта \\
\bottomrule
\end{tabular}
\end{table}

\subsection{Корпуса соответствия}

В корпусе подписанных записей каждый случай отказа несёт точный ожидаемый код и ссылку на положительный контроль, поэтому реализация, которая отказывает всегда, корпус не проходит. Генератор на Go пишет случаи и ожидания. Проверка на Python сама канонизирует байты и проверяет подписи Ed25519 и P-256, вердиктов Go не импортируя. В репозитории лежат 944 именованных случая, а 1688 механических развёрток материализуются по манифесту с хэшами SHA-256. Тест библиотеки генерирует полный корпус, сверяет манифест и требует точного кода и плана для всех 2632 случаев. Мы повторили его на публичной ревизии f76a400 библиотеки (Go~1.26.0, darwin/arm64): 2632 подтеста прошли, код выхода 0. Процедура корпуса задаёт и отрицательный контроль самих инструментов: изменённый байт подписи в положительном случае обязан дать расхождение у генератора и код \texttt{envelope\_signature\_invalid} у проверки на Python.

Для памяти первый вычислитель написан по сводному контракту ядра без доступа к пакету и прогнан на векторах пакета без их правки, второй написан только по пакету, и каждое расхождение разобрано как ошибка пакета, реализации или неоднозначность контракта. На корпусе 0.2.7 оба дают ожидаемые ответы на всех 66 смысловых векторах; в текущем черновике 0.3.0 вычислитель на Go проходит все 77 смысловых векторов.

\subsection{Привилегированные шаги хоста станции}

Владелец выполнил на одной машине Linux два привилегированных запуска, каждый на своей сборке. Первый прошёл режимы применения, проверки и отката для одноразовых примитивов учёток: команды учёток и групп, защита дома, отказ ACL и capabilities, остановка через pidfd с проверкой отсутствия процессов и ограниченный откат дома, учётки и группы. Второй установил набор квалификации: при применении квалификация в режиме загрузки прошла 11 из 11 проверок, отдельная проверка тоже прошла 11 из 11, а допуск станций остался не настроен. Опубликованная запись о результатах указывает, что второй запуск не устанавливает ни кампании перезагрузок, ни отката установки, а проверка открытых inode и пространств имён остаётся неисполненной.

\subsection{Ограниченная модель интеграции}

Модель в \texttt{paper/model/} воспроизводится командой \texttt{make -C paper model}. Она содержит две задачи, до двух версий контракта, до двух эпох аренды и до двух доставок записи и на границе доставки применяет четыре правила: голова равна базе записи (\texttt{cas}), запись несёт текущую эпоху (\texttt{epoch}), аттестация ссылается на текущий контракт (\texttt{contract}) и адресована точному кандидату (\texttt{bind\_I}). При всех правилах перебор в ширину достигает 1090 состояний без нарушения. Без \texttt{cas} запись применяется дважды после 84 просмотренных состояний, без \texttt{epoch} устаревшая эпоха продвигает голову после 82, без \texttt{contract} аттестация ссылается на заменённый контракт после 71, без \texttt{bind\_I} аттестация сделана на иной базе после 490. Эти числа отмечают остановку поиска и безопасность вариантов не сравнивают.

\subsection{Микрозамеры эталонной проверки записей}

Мы измерили шесть публичных вызовов библиотеки проверки на ревизии 9b42338, на неизменённых входах корпуса; реализация эталонная и не оптимизировалась. Замер шёл в гостевой системе KVM на процессоре AMD Ryzen 7 9700X (гостю видны 8 логических CPU, без закрепления), ядро Linux 7.0.0-38-generic, Go 1.26.10 linux/amd64, \texttt{GOMAXPROCS=1}, пять последовательных повторов стандартного бенчмарка Go с \texttt{-benchmem}. Машина была общей: одноминутная средняя нагрузка составляла 1,73 до замера и 1,95 после. Внутри таймера библиотека разбирает и проверяет сырые подписанные байты, и каждая итерация проверяет, что вызов принят; загрузка фикстур, ввод-вывод и подпись в замер не входят.

\begin{table}[t]
\centering\small
\caption{Эталонная проверка записей на вызов. Время дано медианой и размахом [мин.--макс.] пяти повторов; размах не является доверительным интервалом. Память ($10^6$ байт) и число выделений даны медианой.}
\label{tab:bench}
\begin{tabular}{L{7.6cm}rrr}
\toprule
Вызов и вход & Время, мс & Память & Выделений \\
\midrule
Сообщение \texttt{message/1}, 835 байт & 6,04 [5,70--7,22] & 6,30 & 92\,619 \\
Грант, одно звено, 832 байта & 9,10 [7,76--10,53] & 9,00 & 132\,307 \\
Грант, три звена, 2628 байт & 17,99 [15,53--23,75] & 18,91 & 278\,351 \\
Аттестация ключа сообщений, 840 байт & 3,78 [3,71--4,51] & 4,49 & 66\,016 \\
Снимок отзывов на 64 записи, 7110 байт & 2,11 [1,67--2,72] & 1,94 & 28\,497 \\
Журнал из 64 записей и две контрольные точки, 30\,718 байт & 9,21 [7,54--11,23] & 7,53 & 112\,915 \\
\bottomrule
\end{tabular}
\end{table}

Разница одного и трёх звеньев не оценивает влияния глубины цепочки: входы различаются содержимым, и в трёхзвенной цепочке есть подпись P-256. Все 64 записи снимка отзывов вступают в силу позже момента проверки, поэтому строка не измеряет худшего поиска среди действующих отзывов. Вызов журнала проверяет 64 записи и сравнивает вторую контрольную точку на раздвоение, второго суффикса он не проверяет. Проверка одной записи занимает миллисекунды на одном ядре, что мало рядом с ходом модели; накладных расходов служб, ввода-вывода и сети замер не включает.

\subsection{Статистика ревью}

Разработка ведётся методом раздела~\ref{sec:design}: продьюсер публикует неизменяемую ревизию, отдельный агент-ревьювер выносит вердикт с находками по классам серьёзности, и доска не даёт запуску принять собственную ревизию. Мы разобрали файлы исходов ревью на доске задач проекта Swarma Curator по состоянию на 24 октября 2026 года. Вердиктом считалось последнее значение поля \texttt{verdict} в файле, серьёзности считались по JSON находок, записи без разбираемого вердикта учитывались только в раундах. Нашлось 140 задач с записями ревью и 318 записей. Из 104 разобранных вердиктов 55 \emph{принято} и 49 \emph{нужны изменения}; из 35 разобранных первых раундов принято 17. Медиана числа раундов на задачу равна 2, максимум 13. Находок высшего класса серьёзности P0 нет, классов P1 и P2 соответственно 73 и 49. Второй разборщик текстовых вердиктов даёт 58 и 51 при 209 неразобранных записях. Ранние ревью шли в других форматах, другие проекты не учтены, поэтому числа служат нижней границей. Они показывают, что ревьювер регулярно возвращает работу, и ничего не говорят о дефектах, которые он пропустил.

\subsection{Ограничения}

Сквозного прогона станции, где отзыв гранта или квалификации проходит через реестр, брокер, диспетчер и хост сессий и измеряется задержка до завершённого сужения или остановки, нет; граница 300~с и формула $\max(S+A,F)$ остаются свойствами дизайна. Оценки производительности системы целиком нет: замеры относятся к библиотеке на горячих входах. Состязательной оценки против живого развёртывания нет: корпуса проверяют согласие реализаций со спецификацией, а инъекции, захват машины и атаки на узел доверия на работающей установке не испытывались. Нет и измерения пользы роя и памяти на внешних задачах. Привилегированные шаги выполнены на одной машине, а модель интеграции не описывает аренды, остановки и внешних эффектов.

\section{Смежные работы}
\label{sec:related}

\textbf{Полномочия.} Гранты продолжают линию capability-систем \cite{dennis1966,miller2006}: право передаётся вместе со ссылкой и только сужается. Macaroons \cite{macaroons} добавляют ограничения к токену без обращения к издателю, Biscuit \cite{biscuit} даёт публично проверяемые токены с офлайн-ослаблением и языком правил наподобие Datalog, SPKI/SDSI \cite{spki} и UCAN \cite{ucan} строят цепочки делегирования на открытых ключах. Swarma заимствует ослабление и проверку у получателя и добавляет правила уровня платформы: ключ агента не бывает издателем, вызывающий остаётся зависимостью отзыва потомка, отзыв распространяется подписанными снимками с явной границей срока.

\textbf{Прозрачность ключей.} Certificate Transparency \cite{ct,ct2} и CONIKS \cite{coniks} делают раздвоение журнала сертификатов и каталога ключей обнаружимым с помощью деревьев Меркла и аудиторов, Sigstore \cite{sigstore} применяет журнал прозрачности к подписям артефактов. Реестр владельца решает более узкую задачу одного владельца: контрольные точки и свидетель ловят откат и раздвоение в пределах засвидетельствованного, а раздвоение перед разными проверяющими остаётся пределом, который снимает только общее состояние цепочки.

\textbf{Самосертифицирующие идентификаторы.} KERI \cite{keri} вводит журнал ключевых событий и предротацию со следующими ключами за хэшем. Swarma использует ту же идею для ячейки личности, ведёт журнал в реестре владельца до запуска цепочки и выводит ключи агентов иерархически \cite{bip32,slip10}, храня ключи владельца отдельно от мастер-секрета установки.

\textbf{Мессенджеры.} Перекрёстная подпись Matrix \cite{matrix} и MLS \cite{mls} удостоверяют устройства и клиентов и дают групповое сквозное шифрование, то есть отвечают, какое устройство отправило сообщение. Swarma использует перекрёстную подпись как уровень A1 и добавляет поверх носителя подпись точных байтов, гранты и ворота доверия: подлинное устройство права на действие не означает.

\textbf{Изоляция.} gVisor \cite{gvisor}, Firecracker \cite{firecracker} и Kata Containers \cite{kata} изолируют недоверенный код ядром в пространстве пользователя или лёгкими виртуальными машинами. Станция опирается на пользователя ОС, группу ресурсов и файрвол ядра без образа окружения, а более сильные классы исполнения входят через тот же контракт изоляции после отдельной квалификации. Вклад станции составляют объявленные гарантии, их статусы и квалификация, применимые к любому бэкенду.

\textbf{Агентские системы.} ReAct \cite{react} связал рассуждение модели с вызовом инструментов, SWE-agent \cite{sweagent} и SWE-bench \cite{swebench} перенесли это на задачи разработки, AutoGen, MetaGPT и ChatDev \cite{autogen,metagpt,chatdev} организуют несколько ролей и обмен между ними. Параллельные агенты, написавшие компилятор \cite{compiler}, опирались на общий репозиторий и тесты как оракул, а Prove2Me \cite{prove2me} объединяет агентов независимых владельцев для формализации математики, где принятие решает ядро доказательств. Эти работы изучают способности и координацию, а ключи, полномочия и отзыв оставляют среде. Инъекция в промпт против агентов с инструментами описана в \cite{greshake} и измерена в AgentDojo \cite{agentdojo}; Swarma отвечает на неё ограничением прав процесса и детерминированным допуском, а распознавание инъекции служит только сужающей проверкой.

\textbf{Ревью кода моделями.} Автоматизация ревью изучалась на предобученных моделях \cite{tufano2021,codereviewer} и в промышленном исправлении замечаний \cite{frommgen2024}. Исследования моделей в роли судьи показывают систематические смещения, в том числе предпочтение собственного текста \cite{llmjudge,selfpref}. Отсюда правило Swarma: ревьювер вычисляется после фиксации кандидата, а независимость оценивается по административному домену, семейству модели и тулчейну. Журнал протокола следует BFT-протоколам \cite{pbft,hotstuff}, а пределы репутации и множественных личностей описаны в \cite{sybil,sybilproof,eigentrust}.

\section{Заключение и дальнейшая работа}

Swarma переносит решение о допуске из суждения модели в детерминированные проверки у исполнителя эффекта. Захваченный агент не владеет ключом, не расширяет грант делегированием, не получает прав от доставки сообщения и не исполняется на станции с неподтверждёнными гарантиями. Доверие сводится к ядру и root хоста, к узлу доверия под администрированием владельца и к человеку с аппаратным ключом, и у каждой гарантии названы допущения и предел. Имеющиеся свидетельства относятся к спецификациям и отдельным компонентам: корпуса согласуют независимые реализации, привилегированные шаги хоста исполнимы на реальной машине, ограниченная модель показывает необходимость правил интеграции, а эталонная проверка записи стоит миллисекунды.

Дальнейшая работа следует ограничениям раздела~\ref{sec:eval}. Первым нужен сквозной прогон станции на одной машине: ключник, реестр и свидетель на узле доверия, брокер, диспетчер, хост сессий и мост, с измеренной задержкой отзыва гранта и квалификации и с программными нарушителями на настоящих границах. Затем нужны состязательная оценка инъекций против живой установки, сравнение роя с одиночным агентом и с параллельной работой без протокола на одинаковых внешних задачах и калибровка ревьюверов посевом дефектов. Для протокола нужны механически проверяемая спецификация аренды, остановки и внешних эффектов \cite{tla} и векторы соответствия для демона среды. Открытыми остаются формат грантов, страницы снимка отзывов, механизм завершения отзыва, формат аттестации мастер-ключа Matrix и квалификация сервера носителя.

{\footnotesize
\bibliographystyle{plain}
\setlength{\itemsep}{0pt}
\bibliography{references}
}

\end{document}
